\documentclass[10pt,a4paper]{amsart}
\usepackage{amsmath,amssymb,amsthm,graphicx}
% Retained preamble for the Stacks Project chapter snapshot in this directory.
% The chapter's own preamble.tex is a copy of the shared Stacks Project preamble; the block
% below is sliced MECHANICALLY, line for line, from it. Nothing was added, removed or
% redefined; the wrapper class, the xr-hyper/hyperref wiring and the externaldocument list
% are omitted because the assembler supplies the document class.
%
% lines taken from preamble.tex: 3, 5, 9, 11, 12, 156, 157, 158, ... (51 lines total)
%
% package coverage restored from the chapter's own preamble (2026-10-13): xy, xy-2cell

\usepackage[all]{xy}
\xyoption{2cell}
\UseAllTwocells

\newenvironment{reference}{\comment}{\endcomment}
\newenvironment{slogan}{\comment}{\endcomment}
\newenvironment{history}{\comment}{\endcomment}
\theoremstyle{plain}
\newtheorem{theorem}[subsection]{Theorem}
\newtheorem{proposition}[subsection]{Proposition}
\newtheorem{lemma}[subsection]{Lemma}
\theoremstyle{definition}
\newtheorem{definition}[subsection]{Definition}
\newtheorem{example}[subsection]{Example}
\newtheorem{exercise}[subsection]{Exercise}
\newtheorem{situation}[subsection]{Situation}
\theoremstyle{remark}
\newtheorem{remark}[subsection]{Remark}
\newtheorem{remarks}[subsection]{Remarks}
\numberwithin{equation}{subsection}
\def\lim{\mathop{\mathrm{lim}}\nolimits}
\def\colim{\mathop{\mathrm{colim}}\nolimits}
\def\Spec{\mathop{\mathrm{Spec}}}
\def\Hom{\mathop{\mathrm{Hom}}\nolimits}
\def\Ext{\mathop{\mathrm{Ext}}\nolimits}
\def\SheafHom{\mathop{\mathcal{H}\!\mathit{om}}\nolimits}
\def\SheafExt{\mathop{\mathcal{E}\!\mathit{xt}}\nolimits}
\def\Sch{\mathit{Sch}}
\def\Mor{\mathop{\mathrm{Mor}}\nolimits}
\def\Ob{\mathop{\mathrm{Ob}}\nolimits}
\def\Sh{\mathop{\mathit{Sh}}\nolimits}
\def\NL{\mathop{N\!L}\nolimits}
\def\CH{\mathop{\mathrm{CH}}\nolimits}
\def\proetale{{pro\text{-}\acute{e}tale}}
\def\etale{{\acute{e}tale}}
\def\QCoh{\mathit{QCoh}}
\def\Ker{\mathop{\mathrm{Ker}}}
\def\Im{\mathop{\mathrm{Im}}}
\def\Coker{\mathop{\mathrm{Coker}}}
\def\Coim{\mathop{\mathrm{Coim}}}
\DeclareMathSymbol{\boxtimes}{\mathbin}{AMSa}{"02}
\def\QCohstack{\mathcal{QC}\!\mathit{oh}}
\def\Cohstack{\mathcal{C}\!\mathit{oh}}
\def\Spacesstack{\mathcal{S}\!\mathit{paces}}
\def\Quotfunctor{\mathrm{Quot}}
\def\Hilbfunctor{\mathrm{Hilb}}
\def\Curvesstack{\mathcal{C}\!\mathit{urves}}
\def\Polarizedstack{\mathcal{P}\!\mathit{olarized}}
\def\Complexesstack{\mathcal{C}\!\mathit{omplexes}}
\def\Pic{\mathop{\mathrm{Pic}}\nolimits}
\def\Picardstack{\mathcal{P}\!\mathit{ic}}
\def\Picardfunctor{\mathrm{Pic}}
\def\Deformationcategory{\mathcal{D}\!\mathit{ef}}
\usepackage{cleveref}
\crefname{originallemma}{original lemma}{original lemmas}
\Crefname{originallemma}{Original lemma}{Original lemmas}
\makeatletter
\newlabel{algebra-lemma-huber}{{07CF}{source}{}{sourcepointer-1}{}}
\newlabel{algebra-lemma-huber@cref}{{[originallemma][0][]07CF}{[1][0][]source}{}{}{}}
\makeatother
\crefname{originallemma}{original lemma}{original lemmas}
\Crefname{originallemma}{Original lemma}{Original lemmas}
\makeatletter
\newlabel{algebra-lemma-locally-free-semi-local-free}{{02M9}{source}{}{sourcepointer-2}{}}
\newlabel{algebra-lemma-locally-free-semi-local-free@cref}{{[originallemma][0][]02M9}{[1][0][]source}{}{}{}}
\makeatother
\crefname{originallemma}{original lemma}{original lemmas}
\Crefname{originallemma}{Original lemma}{Original lemmas}
\makeatletter
\newlabel{lemma-syntomic-quasi-finite}{{0BWE}{source}{}{sourcepointer-3}{}}
\newlabel{lemma-syntomic-quasi-finite@cref}{{[originallemma][0][]0BWE}{[1][0][]source}{}{}{}}
\makeatother
\title{Stacks, Discriminants: lemma-syntomic-finite}
\author{The Stacks Project Authors}
\date{stacks-project:discriminant.tex:local-snapshot-20261012; GNU FDL 1.2 or later}
\begin{document}
\maketitle
\noindent\textit{Editorial note.} Statement span is lines 2281-2299 of sources/stacks-discriminant/pinned-original-source.tex; proof span is lines 2301-2323, the complete author proof (21 body lines).\par\medskip
\noindent\textit{Editorial note.} Navigation stubs: 3 cited-but-not-included labels are supplied as non-load-bearing pointers with their pinned tags where the snapshot tag index has them.\par\medskip
\noindent\textit{Editorial note.} extra\_packages is empty on purpose: the retained preamble already loads xy with options and a second declaration clashes in the assembler wrapper.\par\medskip
\noindent\textit{Editorial note.} Text is retained verbatim, including the author's own typos and placeholders.\par\medskip
\noindent\textit{Editorial note.} Queue index 24 of queue\_discriminant.json (pb=1098, stmt 2281-2299, proof 2301-2323).\par\medskip
\noindent\textit{Allowed standard prerequisites (editorial).}\begin{itemize}
\item Explicit prerequisite boundary. The cited inputs are prerequisites; the argument in the retained span is not replaced by any cited result.
\end{itemize}\medskip
\ifdefined\hypertarget\hypertarget{sourcepointer-1}{}\fi
\noindent\textit{Source reference.} Original lemma 07CF: sources/stacks-discriminant/pinned-original-source.tex, source lines 2323--2323. Cited by the retained statement or proof but not part of the collected spans; supplied only as a navigation pointer. The argument does not depend on its content.\par\medskip
\ifdefined\hypertarget\hypertarget{sourcepointer-2}{}\fi
\noindent\textit{Source reference.} Original lemma 02M9: sources/stacks-discriminant/pinned-original-source.tex, source lines 2323--2323. Cited by the retained statement or proof but not part of the collected spans; supplied only as a navigation pointer. The argument does not depend on its content.\par\medskip
\ifdefined\hypertarget\hypertarget{sourcepointer-3}{}\fi
\noindent\textit{Source reference.} Original lemma 0BWE: sources/stacks-discriminant/pinned-original-source.tex, source lines 2005--2026. Cited by the retained statement or proof but not part of the collected spans; supplied only as a navigation pointer. The argument does not depend on its content.\par\medskip
\begin{lemma}
\label{lemma-syntomic-finite}
Let $f : Y \to X$ be a morphism of schemes. The following are equivalent
\begin{enumerate}
\item $f$ is finite and syntomic,
\item $f$ is finite, flat, and a local complete intersection morphism,
\item $f$ is finite, flat, locally of finite presentation,
and the fibres of $f$ are local complete intersections,
\item $f$ is finite and for every $x \in X$ there is an
affine open $x \in U = \Spec(A) \subset X$ an integer $n$
and $f_1, \ldots, f_n \in A[x_1, \ldots, x_n]$ such that
$f^{-1}(U)$ is isomorphic to the spectrum of
$A[x_1, \ldots, x_n]/(f_1, \ldots, f_n)$,
\item $f$ is finite, flat, locally of finite presentation,
and $\NL_{X/Y}$ has tor-amplitude in $[-1, 0]$, and
\item $f$ is finite, flat, locally of finite presentation, and
$\NL_{X/Y}$ is perfect of rank $0$ with tor-amplitude in $[-1, 0]$,
\end{enumerate}
\end{lemma}

\begin{proof}
The equivalence of (1), (2), (3), (5), and (6)
and the implication (4) $\Rightarrow$ (1) follow immediately
from Lemma \ref{lemma-syntomic-quasi-finite}. Assume the equivalent conditions
(1), (2), (3), (5), (6) hold.
Choose a point $x \in X$ and an affine open $U = \Spec(A)$
of $x$ in $X$ and say $x$ corresponds to the prime ideal
$\mathfrak p \subset A$. Write $f^{-1}(U) = \Spec(B)$.
Write $B = A[x_1, \ldots, x_n]/I$. Since $\NL_{B/A}$
is perfect of tor-amplitude in $[-1, 0]$ by (6)
we see that $I/I^2$ is a finite locally free $B$-module
of rank $n$. Since $B_\mathfrak p$ is semi-local we see that
$(I/I^2)_\mathfrak p$ is free of rank $n$, see
Algebra, Lemma \ref{algebra-lemma-locally-free-semi-local-free}.
Thus after replacing $A$ by a principal localization at
an element not in $\mathfrak p$ we may assume $I/I^2$
is a free $B$-module of rank $n$.
Thus by Algebra, Lemma \ref{algebra-lemma-huber}
we can find a presentation of $B$ over $A$
with the same number of variables as equations. In other words,
we may assume $B = A[x_1, \ldots, x_n]/(f_1, \ldots, f_n)$.
This proves (4).
\end{proof}

\end{document}
